% Ad Familiares 10.13 --- headnote
% ad-familiares-10-13, 11 May 43 BC, Rome

\textit{Cicero to Plancus, written from Rome around 11 May 43 BC
--- Perseus dateline \textit{Scr. Romae circ. v Id. Mai. a. 711
(43)}. A short covering note for a senatorial decree in
Plancus's honour. News of the victory at Mutina (21 April)
has reached the city; Antony has been driven back across the
Alps; Plancus, governor of Transalpine Gaul, has put his army
in motion and written to the Senate of his loyalty. Cicero,
who has been the architect of the senatorial campaign against
Antony since the previous September, drafts the motion of
thanks himself and reports back to Plancus that it has passed
in a full house with overwhelming support.}

\textit{The second section delivers the political point under
the compliment: Plancus must not slacken. The war is not over
when Antony is beaten; the war is over when Antony is
finished. Cicero reaches for Homer to make the case ---
Ulysses, not Ajax or Achilles, is the \textit{[Greek:
ptolipor\=thion]}, the sacker of cities, because it was
Ulysses who delivered the final stroke. The flattery is
calculated: Plancus is being told he can be the Ulysses of
this war if he closes it now.}
